% Generated by report/make_restated.py from theory/A1_lower_bound.tex: do not edit (edit the note and rerun).
\begin{restated}{Corollary}{cor:sweep}{A1}{1}{initial full gradient}
On the instance \eqref{eq:instance} with $\zeta=\zeta_Q$, and $Q$, $q$ as in Proposition~\ref{prop:minimizer}, let $\theta_0=0$ and let the method be of class (S1) (Definition~\ref{def:span}) preceded by one full gradient at $\theta_0$ (Lemma~\ref{lem:support}(ii)), $t$ denoting the number of i.i.d.\ queries made after this initial full gradient. Then for every $t\ge0$
\begin{equation*}
\tag{\ref*{ineq:expected-sweep}}
\frac{\E\norm{\theta_t-\etastar}^2}{\norm{\etastar}^2}\ \ge\ \frac{q^2\big(1-\frac{1-q^2}n\big)^t-q^{2m}}{1-q^{2m}},
\end{equation*}
and, under \eqref{ineq:dim} and for $\varepsilon<q^2/2$, reaching $\E\norm{\theta_t-\etastar}^2\le\varepsilon\norm{\etastar}^2$ requires
\[
t\ \ge\ \Big(\frac{n(\sqrt Q+1)^2}{4\sqrt Q}-1\Big)\log\frac{q^2}{2\varepsilon}\ \ge\ (n-1)\log\frac{q^2}{2\varepsilon}
\]
i.i.d.\ queries, i.e.\ $G=n+t$ gradient evaluations in total. Compared with \eqref{eq:iidcount}, $\varepsilon$ is replaced by $\varepsilon/q^2$: the initial full gradient reveals the first coordinate of every block, which carries the fraction $1-q^2$ of the energy of each block, and this is worth about $n$ i.i.d.\ queries since $\big(\frac{n(\sqrt Q+1)^2}{4\sqrt Q}-1\big)\cdot2\log\frac1q\le\frac{n(\sqrt Q+1)^2}{\sqrt Q(\sqrt Q-1)}$, which is $n(1+O(Q^{-1/2}))$ for large $Q$ (for $\kappa\le n$, however, $q^2\le0.03$ and the shift $(n-1)\log(1/q^2)$ is a noticeable part of the count).
\end{restated}
